\subsection{Exact computer-assisted certificate format for the arithmetic tables}

Each certificate used in Section~\ref{sec:chebotarev} is a finite identity in
an explicitly specified ring, checked by the hash-pinned verifier of
\Cref{prop:pinned-normalizer-certificate-table}.  The recurrence certificate is
the exact residual certificate \(\mathsf W_q(P_q)=\mathrm{accept}\) of
\Cref{def:exact-kernel-verifier}, together with the scalar-free construction
lemma \Cref{lem:scalar-free-verifier}, the a priori certification theorem
\Cref{thm:apriori-exact-kernel-certification}, and the recurrence theorem
\Cref{prop:kernel-derived-recurrences}.  It starts from the exact
histogram kernel of \Cref{def:exact-histogram-kernel}, checks the exact integer
identity
\[
  R_qP_q(\widetilde A_q)\widetilde v_q=0,
\]
which is the finite Cayley--Hamilton construction of
\(P_q(\widetilde A_q)\widetilde v_q\in Z_2\).  Membership in \(Z_2\) gives the
all-tail recurrence.  The additional row-reduction identities
\(\rk(R_qK_q)=\rk(R_qB_q)=d_q\) form the observable-quotient certificate
\(\mathsf V_q(P_q)=\mathrm{accept}\); they identify the quotient polynomial and
support the Hankel and minimality arguments, but they are not required for the
tail recurrence.
Thus a recurrence row is used only after it has passed through the chain
\[
  \begin{array}{c}
  \text{histogram kernel}\\
  \Downarrow\\
  R_qP_q(\widetilde A_q)\widetilde v_q=0\\
  \Downarrow\ \text{membership in }Z_2\\
  P_q(E)S_q(j)=0\quad(j\ge2).
  \end{array}
\]

The remaining tables have the following exact meanings.

A Smith-normal-form certificate consists of integer unimodular matrices \(U,V\)
with
\[
  U M_q V=\operatorname{diag}(1,\dots,1)
\]
for the relevant basis-change matrix \(M_q\) (with the vacuous \(0\times0\)
case when \(\kappa_{12}=0\)).  A mod-\(p\) factorization certificate consists
of monic factors in \(\FF_p[\lambda]\) whose product is \(P_q\bmod p\), together
with Bezout polynomials certifying
\(\gcd(P_q\bmod p,P_q'\bmod p)=1\).  These data are the artifact fields
\(\texttt{finite\_field\_factorization\_records}\).  A
discriminant-squareclass certificate computes
\(\Disc(P_q)=(-1)^{d_q(d_q-1)/2}\operatorname{Res}(P_q,P_q')\) and then applies
Euler's criterion in the displayed finite fields; the artifact field
\(\texttt{discriminant\_squareclass\_records}\) records the exact
discriminants, residues, and Euler witnesses.  Thus the tables record digest
views of exact algebraic certificates.  The full witness records are the
submission-included JSON fields checked by the displayed verifier command; the
TeX tables are compact, human-readable indices to those exact records.  The
high-precision root table below is the only numerical table, and it is not used
in the proofs.

\medskip
\noindent\textbf{Role in the paper.}
The proofs in Section~\ref{sec:chebotarev} use only the exact polynomials
\(P_q\) determined by the kernel-certified recurrences of
Table~\ref{tab:appendix-recurrences-q9-17} and
\Cref{prop:kernel-derived-recurrences}, the Smith-normal-form certificates
displayed in
\Cref{tab:appendix-principalization-q9-17} together with the recorded
unimodular matrices, and the finite mod-\(p\) factorization records and
discriminant witnesses summarized in
\Cref{tab:galois-certificates-q9-17,tab:disc-squareclass-q12-15}.  The
high-precision root fingerprints recorded in
\Cref{tab:secondary-spectral-q9-17} are included as secondary descriptive data,
but they are not used in the rigorous theorem chain.
The appendix is included to make these inputs explicit and citable inside the
submission itself, with the exact witness data carried by the pinned local
certificate file rather than by the secondary numerical fingerprints.


